%Abstract
\setlength{\headheight}{14pt}
\begin{center}
	{\huge \bf Abstract}
	\line(1,0){430}
\end{center}

Recorded lectures and digital course material are now abundant, but the tools
that help students convert that material into practice and self assessment are
not. Existing study aids either require the student to paste in clean text by
hand, or depend on paid cloud APIs that are impractical in bandwidth limited and
cost-sensitive settings. This report presents \textbf{LectureAssist}, an
end-to-end agentic AI system that turns a raw lecture video or an academic
document into structured study and assessment material without any proprietary
cloud service. The system couples a dual-channel extraction layer scene-adaptive
optical character recognition that distinguishes slides, blackboards, whiteboards
and on-screen text, combined with Whisper-based speech transcription to a
multi-agent generation pipeline that plans, retrieves, generates, validates and
refines five types of quiz questions, all served by locally hosted Gemma~3 models
through the Ollama runtime on a single consumer-grade GPU. A retrieval augmented
chat interface and an adaptive difficulty engine that tracks per topic performance
complete the loop from raw media to personalised assessment. The system was
evaluated on a large scale study in which multiple-choice questions were generated
from a college-level calculus chapter across three difficulty levels, and scored by
an independent, larger judge model on question quality, grounding, diversity and
separately --- \emph{answer correctness}, using a second Answer Generator agent
that re-solves each question without seeing the marked option. The work demonstrates that high quality, personalised educational AI
is achievable on accessible hardware, and that question generators must be judged
on answer correctness and diversity, not on fluency alone.
\clearpage
\setlength{\headheight}{12pt}
